%==============================================================================
% design-caixas.tex — Caixas pedagógicas (tcolorbox) + retrocompatibilidade
%==============================================================================
% Estilo base comum
\tcbset{
  caixapedag/.style={
    enhanced, breakable, arc=2pt, boxrule=0.6pt,
    left=8pt, right=8pt, top=6pt, bottom=6pt,
    fonttitle=\sffamily\bfseries\color{white},
    coltext=black
  }
}

% Fábrica: \novacaixa{ambiente}{cor}{ícone}{Título}
\newcommand{\novacaixa}[4]{%
  \newtcolorbox{#1}{caixapedag, colframe=#2, colback=#2!6,
    colbacktitle=#2, title={#3~#4}}%
}

\novacaixa{objetivos}{corObjetivo}{\faBullseye}{Objetivos de aprendizagem}
% NOTA: abntex2 já define o ambiente `resumo` (resumo/abstract da obra).
% Para não quebrar a retrocompatibilidade (o livro usa \begin{resumo} em
% prefacio/06-resumo.tex e 07-abstract.tex), a caixa pedagógica de síntese
% de capítulo usa o nome `resumocap`.
\novacaixa{resumocap}{corResumo}{\faListUl}{Resumo \& checklist}
\novacaixa{atencao}{corAtencao}{\faExclamationTriangle}{Atenção}
\novacaixa{dica}{corDica}{\faLightbulb}{Dica}
\novacaixa{errocomum}{corErro}{\faBug}{Erro comum}
\novacaixa{pareepense}{corPensar}{\faQuestionCircle}{Pare e pense}
\novacaixa{exercicio}{corExercicio}{\faPencilRuler}{Exercício}
\novacaixa{analogia}{corAnalogia}{\faProjectDiagram}{Analogia}

%------------------------------------------------------------------
% Retrocompatibilidade: comandos legados no novo visual
%------------------------------------------------------------------
\renewcommand{\spec}[1]{%
  \begin{tcolorbox}[caixapedag, colframe=corPrimaria, colback=corPrimaria!6,
    colbacktitle=corPrimaria, title={\faFileContract~SDD — Especificação}]%
  #1
  \end{tcolorbox}}

\renewcommand{\teste}[1]{%
  \begin{tcolorbox}[caixapedag, colframe=corTesteSucesso, colback=corTesteSucesso!6,
    colbacktitle=corTesteSucesso, title={\faCheckCircle~Teste TDD}]%
  #1
  \end{tcolorbox}}

\renewcommand{\testefalha}[1]{%
  \begin{tcolorbox}[caixapedag, colframe=corTesteFalha, colback=corTesteFalha!6,
    colbacktitle=corTesteFalha, title={\faTimesCircle~Teste TDD falhou}]%
  #1
  \end{tcolorbox}}

\renewcommand{\destaque}[1]{%
  \begin{tcolorbox}[caixapedag, colframe=corDestaque, colback=corDestaque!6,
    colbacktitle=corDestaque, title={\faStar~Destaque}]%
  #1
  \end{tcolorbox}}

\renewenvironment{pratica}[1][]{%
  \begin{tcolorbox}[caixapedag, colframe=corSecundaria, colback=corSecundaria!6,
    colbacktitle=corSecundaria, title={\faFlask~Prática: #1}]%
}{%
  \end{tcolorbox}}
